\section{Conclusion}
\label{sec:conclusion}

Sampling-based performance prediction of configurable systems is routinely
evaluated under a cost term equal to the number of measured
configurations. That count-based abstraction is convenient and ubiquitous,
but it has not been empirically validated against actual measurement
time. This paper asks whether sample size $n$ is a calibrated proxy for
empirical training measurement cost, and whether replacing the cost term
changes the scientific conclusions previously drawn about sampling
strategies.

We formalize an empirical training cost $\TrainCost$ that excludes
validation measurements (which the strategy does not control), define the
count-based baseline cleanly, and instantiate composite and Pareto
formulations under both cost models. We empirically test the validity of
$n$ as a proxy across ten configurable systems, including a deeper x264
study over $17$ inputs and two non-functional properties.

The data does not support a single answer. Across seven non-x264
configurable systems, sample size $n$ is a strongly correlated proxy for
$\TrainCost$ ($\rho\geq 0.88$) but is not always calibrated: median
$\mathit{TrainCostPct}/\mathit{SizePct}$ drifts to $0.80$ on
BerkeleyDBC and $0.86$ on lrzip, and per-strategy IQRs reach
$[0.19,\,1.09]$. Decision-level consequences split sharply: on Dune,
HIPAcc, LLVM, Polly the size-as-cost abstraction is essentially safe
($0/7$ rank flips, Kendall $\tau\geq 0.94$); on lrzip it is clearly
unsafe ($5/7$ flips, Pareto Jaccard $0.6$). On x264, the cost-target
scenario alone moves top-$1$ rank agreement from $0.95$ to $0.82$ and
median Pareto Jaccard from $0.43$ to $0.38$ when cost stays as time but
the prediction target switches from time to size. The size-as-cost
abstraction is therefore neither uniformly safe nor uniformly broken; it
is conditionally valid, and the conditions can be diagnosed with a
small pilot. We recommend that configurable-systems sampling studies
report both $\SizeCost$ and $\TrainCost$ alongside accuracy, so the
cost-model choice is visible rather than hidden.

The twelve pre-stated hypotheses summarized in
\Cref{tab:hypothesis-verdicts} make this conditional validity precise:
H1a, H1b, H2a, H2c are supported on every scope where $\TrainCost$ is
non-trivial; H2b is supported only when the calibration drift clears a
$5\%$ practical floor; H3a, H4a, H4b are supported existentially as
soon as any system or x264 context flips, which is common; and the
cross-system H3b correlation between cost CV and rank-flip count
($\rho=0.90$, $p=6.3\!\times\!10^{-3}$) ties the heterogeneity story
of H1 directly to the decision-level instability of H3 and H4.

The Linux cross-version replication
(\StudyLinux family in \Cref{sec:results}) sharpens two points across four
releases (v4.13, v4.15, v4.20, v5.0, v5.04, v5.07, v5.08), and the right way to read it
is strategy-conditional rather than system-wide: \emph{on Linux,
calibration is a property of the strategy, not the system}. First,
the aggregate H2b verdict on Linux \emph{refutes} the literal
hypothesis on every release because the per-strategy ratios average
to $\sim\!\!1.00$, but the per-strategy decomposition shows
\texttt{diversity} miscalibrated at $\sim\!\!2\times$ on all seven
releases (medians $2.28, 2.23, 2.27, 2.06, 2.09, 2.00, 2.02$ (window $[2.00, 2.28]$) --- seven independent
randconfig pools, magnitudes within $\pm 15\%$ of one another, the
most robust positive finding in the paper) while
\texttt{random}/\texttt{balanced}/\texttt{coverage} are calibrated
within $\\pm 0.7\\%$ of $1$ on every release;
\texttt{uncertainty} drifts in pool-dependent direction with
absolute $n$ (cheap at small $n$, expensive at large $n$ on every
release; cross-over $n$ varies with pool geometry). The operational
rule on Linux is therefore not ``always use empirical cost'' nor
``$n$ is a fine proxy'' but ``report empirical $\TrainCost$
whenever the strategy panel includes AL or diversity-aware samplers,
or whenever the level grid lands on tied absolute $n$ values where
MRE is sub-threshold; for cost-blind random/balanced/coverage at
percentage levels that separate cleanly, $n$ remains a usable
proxy.'' Second, the H3b cross-system correlation observed on
\StudyReplication ($\rho\!=\!+0.90$ between cost CV and rank-flip count) does
not extend to the Linux pool. Three of the seven releases share the
same $2/7$ flip count across a $5\times$ range of pool sizes
($19{,}952$ to $92{,}562$) and a $1.27\times$ range of cost CVs
($1.12$ to $1.43$); v5.08 alone flips $6/7$, and the cause is its
specific level grid landing on tied absolute $n$'s where MRE
differences are sub-threshold, not pool size or heterogeneity. The Linux quartet therefore replicates the structure of
the H1--H4 verdicts while exhibiting a regime that the
synthesis-based panel cannot produce. The \StudyEnergy generated-sampling
study (\StudyEnergy family in \Cref{sec:results}) sharpens the same conclusion in
the opposite direction: top-$1$ rank flips reach $7/7$ on three of
brotli's four cost-target scenarios and Pareto fronts differ in
$28/28$ scenarios, while PostgreSQL\,11.2 ($\mathrm{CV}\!=\!0.06$)
keeps both fronts almost identical. The size-as-cost abstraction
spans the full failure-to-safety range within a single protocol.

\paragraph*{Implications for practitioners.}
A practitioner running a sampling-based performance-prediction study should
\emph{stop reporting sample size $n$ as the only effort axis}. Instead:
\begin{itemize}
    \item Audit per-configuration cost heterogeneity on a uniform-random
    pilot of $\sim\!\!2$--$5\%$ of $|C|$. If $\mathrm{CV}\!>\!0.5$ or
    $p_{90}/p_{10}\!>\!4$, the size-as-cost abstraction is unsafe on
    that system.
    \item Compute the calibration ratio
    $\mathit{TrainCostPct}/\mathit{SizePct}$ \emph{per strategy} on the
    pilot, not just an aggregate median. Strategy-specific bias
    survives even when aggregate medians look calibrated.
    \item Be especially careful with diversity-aware and active-learning
    samplers. On Linux, \texttt{diversity} pays $\sim\!\!2\times$ its
    count budget on every release; \texttt{uncertainty} drifts in a
    pool-size-dependent direction. On x264, $t$-wise pays
    $\sim\!\!1.10\times$ random's cost and Henard $\sim\!\!0.75\times$.
    Cost-blind samplers (\texttt{random}, \texttt{balanced},
    \texttt{coverage}) are usually fine.
    \item Distinguish the cost dimension from the prediction target. The
    asymmetric case (target $\neq$ cost, \eg predict size while paying
    time-based cost) is the case where the abstraction breaks most
    often: x264's worst scenario at top-$1$ agreement $0.82$ is exactly
    that one. If wall-clock is the binding constraint, cost is time;
    if storage is the binding constraint, cost is size; treat them as
    independent choices.
    \item Report both $\SizeCost$ and $\TrainCost$ alongside any accuracy
    metric. The cost-model choice is then visible to future readers
    rather than hidden inside an unspoken convention.
\end{itemize}

\paragraph*{Implications for researchers.}
The configurable-systems performance-prediction literature has used
count-based effort as the default for over a decade. The data here is
neither a wholesale refutation nor a wholesale endorsement of that
convention; it is a diagnosis of where it fails and where it does
not. Concretely:
\begin{itemize}
    \item Past sampling-cost results should be re-examined when the
    system has high cost CV, when the prediction target differs from
    the acquisition cost dimension, or when the strategy panel
    includes diversity-aware or active-learning samplers. On
    LLVM, Polly, Dune, HIPAcc, and PostgreSQL the size-as-cost
    abstraction stands; on lrzip, BerkeleyDBC, jump3r, brotli, and
    Linux's \texttt{diversity}/\texttt{uncertainty} paths, it does
    not, and previously published rankings deserve restatement under
    $\TrainCost$.
    \item Aggregating sampling-cost results across protocols
    (synthesis-based, pool-based, generated-sampling) is unsafe.
    Each protocol exposes a different strategy panel and a different
    heterogeneity regime, and our cross-system H3b correlation
    ($\rho\!=\!0.90$ between cost CV and rank-flip count on \StudyReplication)
    holds within a fixed protocol but does not transfer to Linux. The
    discipline should report per-protocol verdicts and resist the
    urge to summarize across them.
    \item ML active learning has known for nearly two decades that
    fewer labeled instances do not necessarily imply lower total
    cost when instance costs vary
    (Settles--Craven--Friedland)~\red{[bib]}. The configurable-systems
    sampling literature has not absorbed this. Adopting the
    heterogeneous-cost framing in our community is overdue, and the
    pilot-based diagnostic above is the lightweight step in that
    direction.
    \item Open perspectives. (a) Replicating the findings under cost
    dimensions beyond time, size, and energy --- bandwidth, carbon,
    monetary cloud cost --- is a clean follow-up; the formulation
    accommodates them unchanged. (b) Cross-domain replication on
    real-runtime ML/HPO benchmarks (LCBench, NAS-Bench-201) would
    test whether the configurable-systems story generalizes; we
    deprioritized those here to keep the empirical claim anchored,
    but they are a natural next step. (c) Closing the loop with
    cost-aware samplers that read the calibration ratio at small
    pilot $n$ and adapt their selection accordingly is the
    \emph{designer-side} version of the diagnostic; the diagnostic
    plus the empirical calibration table in
    \Cref{tab:systems-and-cost} provides the prior such a sampler
    would need.
    \item Most importantly, the discipline should treat the cost axis
    as a first-class methodological variable, not a convention. Two
    studies disagree on which strategy is best because they used
    different cost axes is a result, not a bug; it should be reported
    that way.
\end{itemize}

The artifact is available at \red{anonymized URL}.
